\begin{table}[htbp]
\centering
\caption{Event-study estimates of the BIP--White majority gap in mental well-being by sex}
\label{tab:event_study_bip_bysex}
\small
\begin{tabular}{lcccc}
\toprule
Period & Male coefficient & Male p-value & Female coefficient & Female p-value \\
\midrule
2019 (base) & 0.00 &  & 0.00 &  \\
2016 & 0.00 & 1.00 & 0.51* & 0.10 \\
2017 & -0.37 & 0.23 & 0.38 & 0.19 \\
2018 & 0.14 & 0.58 & 0.24 & 0.36 \\
2020/04 & -1.84*** & 0.00 & -0.49 & 0.13 \\
2020/05 & -0.78** & 0.02 & -0.14 & 0.67 \\
2020/06 & -1.21*** & 0.00 & -0.78** & 0.02 \\
2020/07 & -1.07*** & 0.00 & -0.46 & 0.18 \\
2020/09 & -0.46 & 0.14 & -0.10 & 0.78 \\
2020/11 & -0.43 & 0.22 & -0.34 & 0.32 \\
2021/01 & -0.79** & 0.03 & 0.23 & 0.54 \\
2021/03 & -0.77** & 0.02 & 0.02 & 0.96 \\
2021/09 & -0.75** & 0.03 & -0.11 & 0.71 \\
Joint pre-trend &  & 0.38 &  & 0.39 \\
\bottomrule
\end{tabular}

\begin{flushleft}
\footnotesize
Notes: The table reports event-study interaction coefficients for the BIP--White majority gap in mental well-being, estimated separately by sex. The male coefficients compare BIP men with White-majority men, while the female coefficients compare BIP women with White-majority women. The 2019 calendar year is omitted as the baseline period. The outcome is the April-adjusted and inverted GHQ-12 Likert score, so higher values indicate better mental well-being. A negative coefficient indicates that BIP respondents experienced a relative deterioration in mental well-being compared with the White majority within the same sex group. Standard errors are clustered at the individual level. *, **, and *** denote significance at the 10\%, 5\%, and 1\% levels, respectively. The row ``Joint pre-trend'' reports the p-value from the joint test of whether all pre-pandemic interaction coefficients are equal to zero.
\end{flushleft}
\end{table}